\section{Experimental research}

This chapter describes the experiments conducted to evaluate the proposed system.

\subsection{Experimental setup}

The experiments were performed on the following hardware:
\begin{itemize}
  \item CPU: 8-core processor, 3.2~GHz;
  \item RAM: 32~GB DDR4;
  \item Storage: 512~GB NVMe SSD;
  \item OS: Ubuntu 24.04 LTS.
\end{itemize}

\subsection{Test datasets}

Three datasets of varying sizes were used for evaluation:

\begin{table}[H]
    \caption{Test datasets} \label{tab:datasets}
    \begin{tabular}{|l|c|c|c|}
    \hline \textbf{Dataset} & \textbf{Records} & \textbf{Size (MB)} & \textbf{Description} \\
    \hline Small   & 1\,000    & 5   & Unit testing \\
    \hline Medium  & 100\,000  & 450 & Integration testing \\
    \hline Large   & 1\,000\,000 & 4\,200 & Stress testing \\
    \hline
    \end{tabular}
\end{table}

\subsection{Performance results}

The system was benchmarked against two baseline implementations.
Table~\ref{tab:performance} shows the results.

\begin{table}[H]
    \caption{Performance comparison (processing time in seconds)} \label{tab:performance}
    \begin{tabular}{|l|c|c|c|}
    \hline \textbf{Method} & \textbf{Small} & \textbf{Medium} & \textbf{Large} \\
    \hline Baseline A       & 0.8  & 45.2  & 512.0 \\
    \hline Baseline B       & 0.5  & 38.7  & 478.3 \\
    \hline Proposed approach & 0.3  & 22.1  & 287.5 \\
    \hline
    \end{tabular}
\end{table}

The proposed approach demonstrates a \textbf{40--44\%} improvement in processing time compared to Baseline~A, and a \textbf{37--40\%} improvement over Baseline~B.

\subsection{Accuracy results}

Accuracy was measured using standard metrics:

\begin{equation} \label{eq:accuracy}
    \text{Accuracy} = \frac{TP + TN}{TP + TN + FP + FN} \times 100\%
\end{equation}

where $TP$, $TN$, $FP$, $FN$ denote true positives, true negatives, false positives, and false negatives respectively.

The proposed approach achieved an accuracy of \textbf{94.2\%}, compared to 89.5\% for Baseline~A and 91.1\% for Baseline~B.

\subsection{Scalability analysis}

The system scales linearly with dataset size, as evidenced by the following relationship:

\begin{equation} \label{eq:scalability}
    T(n) = a \cdot n + b
\end{equation}

where $T(n)$ is the processing time, $n$ is the number of records, and $a$, $b$ are constants determined through regression analysis ($a = 2.87 \times 10^{-4}$, $b = 0.12$).

\clearpage
